\section*{Two-Component Edition}
\addcontentsline{toc}{section}{Two-Component Edition}

This edition preserves the arguments, normalizations, equation numbering,
and mostly-plus metric conventions of the original text, while expressing
all four-dimensional fermions and superspace formulas in dotted and
undotted two-component notation.  An undotted index
\(\alpha,\beta,\ldots\) labels a \((1/2,0)\) spinor, and a dotted index
\(\dot\alpha,\dot\beta,\ldots\) labels a \((0,1/2)\) spinor.  Thus a
four-component Majorana field in the original notation is represented here
by a Weyl field \(\psi_\alpha\) together with its complex conjugate
\(\bar\psi_{\dot\alpha}\).

We retain
\[
  \eta_{\mu\nu}=\operatorname{diag}(-1,+1,+1,+1)
\]
and define
\[
  \sigma^\mu_{\alpha\dot\alpha}
  =(\mathbf{1},\sigma^1,\sigma^2,\sigma^3),
  \qquad
  \bar\sigma^{\mu\,\dot\alpha\alpha}
  =(\mathbf{1},-\sigma^1,-\sigma^2,-\sigma^3).
\]
It follows that
\[
  \sigma^\mu\bar\sigma^\nu+\sigma^\nu\bar\sigma^\mu
  =-2\eta^{\mu\nu}\mathbf{1},
  \qquad
  \bar\sigma^\mu\sigma^\nu+\bar\sigma^\nu\sigma^\mu
  =-2\eta^{\mu\nu}\mathbf{1}.
\]
These choices reproduce the Dirac matrices used in the original text by
\[
  \gamma^\mu=-\ii
  \begin{pmatrix}
    0&\sigma^\mu\\
    \bar\sigma^\mu&0
  \end{pmatrix},
\]
which fixes all relative phases in the translations below.
We use the Lorentz generators
\[
  \sigma^{\mu\nu}
  =\frac14(\sigma^\mu\bar\sigma^\nu-\sigma^\nu\bar\sigma^\mu),
  \qquad
  \bar\sigma^{\mu\nu}
  =\frac14(\bar\sigma^\mu\sigma^\nu-\bar\sigma^\nu\sigma^\mu).
\]

Spinor indices are raised and lowered with
\[
  \epsilon_{12}=+1,\qquad
  \epsilon^{12}=-1,\qquad
  \psi^\alpha=\epsilon^{\alpha\beta}\psi_\beta,\qquad
  \psi_\alpha=\epsilon_{\alpha\beta}\psi^\beta,
\]
and likewise for dotted indices.  We abbreviate
\(\psi\chi=\psi^\alpha\chi_\alpha\) and
\(\bar\psi\bar\chi
=\bar\psi_{\dot\alpha}\bar\chi^{\dot\alpha}\).

The supersymmetry and superspace operators are written
\[
  Q_\alpha,\quad\bar Q_{\dot\alpha},\qquad
  \mathcal Q_\alpha,\quad\bar{\mathcal Q}_{\dot\alpha},\qquad
  D_\alpha,\quad\bar D_{\dot\alpha}.
\]
All Grassmann derivatives are left derivatives, as in the original text.
A left-chiral superfield obeys
\(\bar D_{\dot\alpha}\Phi=0\), while its Hermitian conjugate obeys
\(D_\alpha\Phi^\dagger=0\).
We bind the squared operators in the orders
\(D^2\equiv D^\alpha D_\alpha\) and
\(\bar D^2\equiv\bar D_{\dot\alpha}\bar D^{\dot\alpha}\).  With the
left-derivative and Berezin-measure conventions of this edition,
\(D^2(\theta^2)=\bar D^2(\bar\theta^2)=-4\),
\((\bar D^2S)^\dagger=D^2S^\dagger\), and
\([\bar D^2h]_{\mathcal F}=-2[h]_D\).  These identities fix the
cross-chapter signs of chiral projections and superspace field equations.

The representation-theory derivations in Chapter 25 that were already
two-component retain Weinberg's temporary Roman component label \(a\).
Equation~\eqref{eq:25.2.34} identifies those components with
\(Q_\alpha\) and \(\bar Q_{\dot\alpha}\); the later chapters use dotted and
undotted indices directly.

Chapter 32 concerns spinors in arbitrary spacetime dimension and therefore
retains the dimension-dependent Clifford notation needed there.  Its
strictly four-dimensional references follow the conventions on this page;
classes such as \(d=4\pmod 8\), which also contain higher dimensions, remain
in the general Clifford notation.
